\chapter{Contesto sperimentale e background fisico}
\label{chap:background}

\paragraph{Riassunto del capitolo.} Il capitolo fornisce il minimo
indispensabile di contesto fisico per leggere il resto del rapporto: il
significato di \emph{Aging}, \emph{Temperatura} e \emph{\acrshort{soc}}
nel nostro dataset, e l'interpretazione della curva di Nyquist come
input \acrshort{ml}. Il dettaglio del banco di misura e del modello a
circuito equivalente non fa parte del nostro lavoro: lo trattiamo come
black-box fornito dal Politecnico di Milano.

\section{Origine del dataset}
\label{sec:experimental-setup}

Il dataset proviene da una campagna di misure condotta presso il
Politecnico di Milano ~\cite{polimi2025battery}. La
caratterizzazione è stata effettuata con tecnica
\acrshort{geis} (\emph{Galvanostatic Electrochemical Impedance
Spectroscopy}) nel range $\sim$\SI{0.1}{\hertz}--\SI{10}{\kilo\hertz},
su una pouch \acrshort{licoo2} commerciale (\cref{tab:cell-specs}), a
varie temperature e stati di carica.

\textbf{Scope del nostro lavoro.} Trattiamo i dati come forniti -- un
insieme di curve $(\mathrm{Re}(Z), \mathrm{Im}(Z))$ campionate -- e non
entriamo nel merito di strumentazione, protocollo di misura o
identificazione dei parametri del modello a circuito equivalente, che
sono di competenza del gruppo di ricerca del Politecnico.

\begin{table}[H]
  \centering
  \caption{Caratteristiche elettriche principali della cella sotto test.}
  \label{tab:cell-specs}
  \begin{tabular}{ll}
    \toprule
    \textbf{Parametro} & \textbf{Valore} \\
    \midrule
    Tipo                                  & Pouch \acrshort{licoo2}, anodo grafite \\
    Capacità nominale                     & \SI{10}{\ampere\hour} \\
    Tensione massima di carica            & \SI{4.2}{\volt} \\
    Tensione minima di scarica            & \SI{2.75}{\volt} \\
    Dimensioni $L \times W \times H$      & $152 \times 72 \times 9$~mm \\
    \bottomrule
  \end{tabular}
\end{table}

\subsection{Mappatura dataset \texorpdfstring{$\to$}{->} paper}

Le variabili provengono dal protocollo
di misura del Politecnico:

\begin{itemize}
  \item \texttt{Aging} $\in \{0,1,2,3,4\}$ -- stadio di invecchiamento
        accelerato della cella; il paper di riferimento studia un
        singolo stato di aging, mentre il dataset esteso che abbiamo
        ricevuto copre cinque livelli;
  \item \texttt{Temperature} $\in \{20.0, 22.5, 25.0, 27.5, 30.0, 35.0,
        40.0, 47.5\}$\,\si{\celsius} -- otto livelli, che campionano
        più finemente la regione $20$--$30$\,°C e si spingono fino a
        $47.5$\,°C;
  \item \texttt{SOC} $\in \{0,1,2,3,4\}$ -- cinque stati di carica
        equispaziati al $25\,\%$ ($0\,\%$, $25\,\%$, $50\,\%$, $75\,\%$,
        $100\,\%$);
  \item \texttt{Frequency} -- 49 punti log-spaziati nel range
        $\sim$\SI{0.1}{\hertz}--\SI{10}{\kilo\hertz};
  \item \texttt{Z\_real}, \texttt{Z\_imag} -- componenti dell'impedenza
        complessa $Z = \mathrm{Re}(Z) + j\,\mathrm{Im}(Z)$, riportate in
        \milliohm{} per leggibilità.
\end{itemize}

\section{La curva di Nyquist come input ML}

Il diagramma di Nyquist riporta $\mathrm{Re}(Z)$ in ascissa e
$-\mathrm{Im}(Z)$ in ordinata, una rappresentazione standard in
spettroscopia di impedenza. Sui dati che usiamo si distinguono due
zone qualitative:
\begin{itemize}
  \item un \emph{semicerchio} alle medie frequenze;
  \item una \emph{coda diffusiva} alle basse frequenze.
\end{itemize}

Le due zone cambiano forma con \emph{Aging}, \emph{Temperatura} e
\emph{\acrshort{soc}}: è esattamente questa variazione che i nostri
modelli imparano a sfruttare. \textbf{Per il nostro lavoro non è
necessario fittare un modello a circuito equivalente né identificare
parametri elettrochimici}: trattiamo le curve come segnali numerici e
le riassumiamo in poche \emph{feature} (vedi
\cref{tab:regression-features,tab:classification-features}).

\section{Effetto di temperatura, SOC e Aging sulle feature ML}

Tre osservazioni qualitative motivano direttamente le scelte di
feature engineering del \cref{chap:methodology}:
\begin{enumerate}[leftmargin=*]
  \item \textbf{Temperatura.} La dipendenza dalla temperatura non è
        lineare; le costanti di tempo dei processi reversibili scalano
        in modo Arrhenius. Pragmaticamente: la feature
        $1 / (T + \SI{273.15}{\kelvin})$ è strutturalmente più
        informativa di $T$ in scala lineare.
  \item \textbf{\acrshort{soc}.} Cambia l'ampiezza del semicerchio e
        l'orientamento della coda diffusiva, ma in modo regolare; basta
        includerlo come feature, eventualmente in interazione con
        $T$ e $Aging$.
  \item \textbf{Aging.} È la dimensione che fa più variare la
        \emph{forma} della curva. Per Task 1 e Task 3 lo usiamo come
        \emph{input}; per Task 2 è il \emph{target} e va escluso dalle
        feature.
\end{enumerate}

A parità di $(Aging, T, \acrshort{soc})$ resta una piccola dispersione
fisiologica delle misure dell'ordine di
$\mathcal{O}(\SI{1}{\milliohm})$, che fissa un \emph{floor} di errore
irriducibile contro il quale nessun modello può andare oltre.
